\section{Random sampling of a viable subspace}

Let $M$ and $F$ be finite-dimensional complex inner product spaces, where $F$
stands for $\mathcal H_L\otimes\mathcal H_R$ and $M$ for $\mathcal H_M$ in
the setting of~\cite[lines 532--547]{Arad2016RigorousRG}; viability is read
up to the associator and the exchange of the first two factors.

\begin{theorem}[Net step]
    \label{thm:viable_random_sampling_net}
    \lean{Submodule.isCloseTo_of_net}
    \leanok
    \uses{thm:subspace_close_iff}
    Let $S$ and $T$ be subspaces of a finite-dimensional inner product space
    and let $c\ge0$. If every unit vector $t\in T$ lies within distance
    $\sqrt c$ of some $v\in T$ with $\lVert P_Sv\rVert^2\ge4c$, then
    $\lVert P_St\rVert^2\ge c\lVert t\rVert^2$ for all $t\in T$, that is, $S$
    is $(1-c)$-close to $T$.
\end{theorem}

\begin{proof}
    \leanok
    For a unit vector $u\in T$ choose $v$ as in the hypothesis. Then
    $\lVert P_Su\rVert\ge\lVert P_Sv\rVert-\lVert u-v\rVert
        \ge2\sqrt c-\sqrt c=\sqrt c$. A general nonzero $t\in T$ is
    $\lVert t\rVert$ times a unit vector of $T$.
\end{proof}

\begin{definition}[Random subspace]
    \label{def:viable_random_subspace}
    \lean{Submodule.randomSubspace, Submodule.finrank_randomSubspace}
    \leanok
    Let $W\subseteq M$ have an orthonormal basis indexed by a finite set
    $\iota$, and let $A\subseteq\iota$ have $s$ elements. For a unitary $U$
    on $\mathbb C^\iota$, the subspace $W'_U\subseteq W$ is the image under
    the basis isometry $\mathbb C^\iota\to W$ of $U$ applied to the span of
    the standard basis vectors indexed by $A$; it has dimension $s$. When $U$ is distributed by the
    Haar measure, $W'_U$ is a uniformly random $s$-dimensional subspace of
    $W$.
\end{definition}

\begin{theorem}[One vector]
    \label{thm:viable_random_sampling_one_vector}
    \lean{Submodule.unitaryHaar_norm_sq_starProjection_randomSubspace_le}
    \leanok
    \uses{def:viable_random_subspace, thm:haar_projection_law,
        thm:complex_gaussian_ratio_lower_tail}
    Let $q=\lvert\iota\rvert$. For every unit vector $f\in W$, the Haar
    probability that $\lVert P_{W'_U}f\rVert^2\le s/(2q)$ is at most
    $e^{-s/16}$.
\end{theorem}

\begin{proof}
    \leanok
    The basis isometry carries projections onto subspaces of
    $\mathbb C^\iota$ to projections onto their images, so the event is the
    event that the projection of a unit vector of $\mathbb C^\iota$ onto
    $U$ applied to the coordinate subspace of $A$ has squared norm at most
    $s/(2q)$. Its probability is the probability that
    $\sum_{i\in A}\lvert g_i\rvert^2\le\frac{s}{2q}\sum_i\lvert g_i\rvert^2$
    for independent standard complex Gaussians $g_i$, which is at most
    $e^{-s/16}$.
\end{proof}

\begin{theorem}[Random sampling]
    \label{thm:viable_random_sampling}
    \lean{Submodule.unitaryHaar_not_isViable_randomSubspace_le}
    \leanok
    \uses{thm:viable_random_sampling_net, thm:viable_random_sampling_one_vector,
        thm:tensor_weights, thm:ball_net, def:subspace_viable}
    Let $T\subseteq M\otimes F$ have dimension $r$, let $W\subseteq M$ be
    $\delta$-viable for $T$ with $\delta<1$, and let $W'_U$ be as above with
    $q=\dim W$. Put
    \[
        \delta'=\frac{(1-\delta)s}{8q},\qquad\nu=\sqrt{\delta'},\qquad
        \eta=\Big(1+\frac2\nu\Big)^{2r}qe^{-s/16}.
    \]
    Then the Haar probability that $W'_U$ is not $(1-\delta')$-viable for
    $T$ is at most $\eta$. This is Lemma~l:rs
    of~\cite[lines 532--547]{Arad2016RigorousRG} with the failure
    probability corrected: the printed value
    $(1+4\sqrt{q/((1-\delta)s)})^rqe^{-s/16}$ undercounts the net of the
    unit sphere of $T$, a real space of dimension
    $2r$~\cite{gap:alvv16_random_sampling_net_count}.
\end{theorem}

\begin{proof}
    \leanok
    If $\delta'=0$ every subspace is $1$-viable. Otherwise choose a
    $\nu$-net $C$ of the unit sphere of $T$ with at most $(1+2/\nu)^{2r}$
    points. For $v\in C$ let $w$ be the projection of $v$ onto $W\otimes F$,
    so $\lVert w\rVert^2\ge1-\delta$ and
    $P_{W'_U\otimes F}v=P_{W'_U\otimes F}w$. Write
    $\lVert P_{W'_U\otimes F}w\rVert^2=\sum_k\mu_k\lVert P_{W'_U}f_k\rVert^2$
    along the weights of $w$. Outside the union of the events
    $\lVert P_{W'_U}f_k\rVert^2\le s/(2q)$, over $v\in C$ and the $q$ basis
    vectors $f_k$, every net point satisfies
    $\lVert P_{W'_U\otimes F}v\rVert^2\ge\frac{s}{2q}(1-\delta)=4\delta'$,
    and the net step gives $(1-\delta')$-viability. The union has
    probability at most $\lvert C\rvert qe^{-s/16}\le\eta$.
\end{proof}

\begin{theorem}[Existence of a viable sample]
    \label{thm:viable_random_sampling_exists}
    \lean{Submodule.exists_isViable_randomSubspace}
    \leanok
    \uses{thm:viable_random_sampling}
    In the setting of the random sampling theorem, if $\eta<1$ then some
    unitary $U$ makes $W'_U$ a $(1-\delta')$-viable subspace for $T$ of
    dimension $s$.
\end{theorem}

\begin{proof}
    \leanok
    Otherwise the failure event is the whole unitary group, of probability
    one, which exceeds $\eta$.
\end{proof}
